\subsection{HOW TO HAVE COMBAT}

\begin{enumerate}
  \item A \textit{possible} Combat situation occurs when enemy units are in each other's Zones of Control. Units that begin their porition of the Turn in enemy Zones of Control, or which move into them are not necessarily going to have Combat. The Non-Phasing Player can elect to avoid Combat, if unwanted.

  \item After the Phasing Player has finished moving his units, the players then consider each possible Combat situation. Each Combat situation is considered in any order determined by the Phasing Player.

  \item The Non-Phasing Player must now decide which of his units in possible combat situations may withdraw immediately and which may not.

  \begin{enumerate}[label=\alph*)]
    \item Stacks containing any of the Non-Phasing Player's infantry and/or artillery units may not withdraw voluntarily until after the first round of Combat.

    \item All Stacks non containing any of the Non-Phasing Player's infantry and/or artillery units may withdraw before any Combat occurs.
  \end{enumerate}

  (NOTE: Infantry and/or artillery units being attacked must always stay in the present hex and accept the first round of Combat. They are free to withdraw before the second or any subsequent rounds of Combat.)

  \item The Non-Phasing Player must now decide what he wants to do in each possible Combat situation. he has two possible options:

  \begin{enumerate}[label=\alph*)]
    \item To remain in his present hex and accept Combat; or
    \item To withdraw one hex to avoid Combat.
  \end{enumerate}

  The Non-Phasing Player must make this decision for each of his units or Stacks of units that lies in an Enemy Zone of Control.

  \item Combat now takes place between any enemy forces still in each others' Zones of Control. Combat continues until no further Combat is possible (this hapepns when the Phasing Player's units no longer have enough Movement Allowance Factors left for another Round of Combat), or until Combat results (retreats) force all units out of each other's Zones of Control.

  \item Once all possible Combat situations have ended, the Phasing Player may move any units that were in possible Combat situations, are no longer in enemy Zones of Control, and still have enough Movement Allowance Factors left to do so. They may lead to more possible Combat situations, withdrawals and/or Combat, still more movement, etc, until the Phasing Player declares that he is finished iwth his Phase of the Turn.

  \textbf{NOTE:} It is possible for the Phasing Player to withdraw from Combat, then using remaining Movement Allowance Factors for further movement (which could result in more Combat).
\end{enumerate}

\subsubsection[GENERAL RULES OF COMBAT]{\mbox{GENERAL RULES OF COMBAT}}

\begin{enumerate}
  \item Combat is fought in Rounds, each Round representing several hours of real fighting. An individual Combat situation may consist of one or more Rounds of Combat.

  \item Each Round of Combat that a unit participates in costs an expenditure of the unit's Movement Allowance Factor - when the Phasing Player's unit has expended all of its Movement Allowance Factor, it can enter no further Rounds of Combat.

  \item Each Round of Combat is divided into two segments; the Phasing Player's segment, and the Non-Phasing Player's segment. During his segment of a Round of Combat, a player may attack any adjacent enemy hex, or adopt other options covered in the Combat Procedure rules.

  \item The Non-Phasing Player's units need not worry about the expenditure of Movement Allowance Factors by Combat. He can continue in action as long as the Phasing Player keeps initiating Rounds of Combat from an adjacent hex. The Non-Phasing Player's units cannot engage in Combat against any of the Phasing Player's units that no longer have the Movement Allowance Factors to start a Round of Combat.

  \item As long as he still has enough Movement Allowance Factors to do so, the Phasing Player must start new Rounds of Combat against any enemy units in an adjacent hex. (NOTE: According to COMBAT PROCEDURES, Rule 3, Section C the player need not actually attack during his segment of Combat, but merely initiate subsequent rounds of combat.)

  \item An enemy occupied hex may be attacked from any or all of the hexes in its Zone of Contro that ocntain the attacker's units.

  \item When you attack a hex, you must attack all units in that hex. When attacked, you must use all your units in the hex for Combat.

  \item The Reserve section of the Stack Boxes plays no part in the Basic Game.

  \item No unit may attack more than once during its segment of the Round of Combat, and no unit may be attacked more than once during the enemy's segment of the Round of Combat.

  \item The units in one hex may attack only one enemy hex during a segment of a Round of Combat.

  \item It is possible for the units in a hex to attack during their own segment of a Round of Combat, and for them to also be attacked during their opponent's segment of that same Round of Combat.

  \item The player taking the casualties has his choice of exactly which types of Combat Factors to remove, unless specified by the COMBAT EFFECTS TABLE.

  \item When attacking, you are under no obligation to attack all enemy held hexes who Zone of Control you are adjacent to.

  \item In no case will an attack against one enemy hex be resolved by more than one die roll on the COMBAT EFFECTS TABLE, reglardless of the number of hexes the attack is made from.

  \item Since units in Battle Formation cannot cross a stream or river hexside, or a bridge, the also cannot attack across these types of terrain. A Non-Phasing Player's units in Battle Formation could defend such a hexside, but would lose their "attack" option during the segment of the Round of Combat. The Phasing Player is not allowed to enter an enemy Zone of Control in Battle Formation if the enemy hex is separate from his by a stream or river hexside.

  \item In the case of withdrawals and retreats, if they are made into the Zones of Control of enemy units, a new First Round of Combat must be initiated if the Phasing Player's units still have enough Movement Allowance Factors to do so. Note that a unit that is retreated into an enemy Zone of Control for further Combat is treated as being in Column Formation - it suffers no penalties for having been forced to retreat.

  \item A unit or Stack of units, both for the Phasing and Non-Phasing Player, is always allowed to occupy (enter) a hex from which an enemy force has just been forced to retreat a result of an attack (you not do not get to occupy the hex as a result of forcing an attacker to retreat). This does not have to be done, but can be made as soon as the retreat takes place. The Phasing Player must use Movement Allowance Factors for this movement, if they are available; but the occupation is still allowed, even if all Movement Allowance Factors have been expendend. You are allowed to occupy, even if a Formation change is necessary for the advance, and even if it is into an enemy Zone of Control.
\end{enumerate}